\section{A finite exterior-power gate for weight--monodromy}
\label{sec:finite-weight-gate}

Scholze (2012), arXiv:1111.4914, Lemma 9.5, printed p.~49,
records Deligne's tensor-power criterion: purity of the monodromy
filtration is equivalent to upper weight bounds on the inertia
invariants of every tensor power and every dual tensor power.
The following finite exterior-power formulation is proved here in
full. It reduces that infinite collection of representations; it
does not assert the needed bounds for an arbitrary geometric
representation. No historical novelty claim is made.

\paragraph{Finding BR-WM-02 (finite exterior-power criterion).}
Let $V$ be a complex vector space of dimension $d\geq1$, let
$q>1$, let $w\in\mathbb R$, and let $F$ be invertible and $N$
nilpotent with $FNF^{-1}=q^{-1}N$. Write $M$ for the monodromy
filtration centered at zero. On exterior powers, $N$ denotes the
induced derivation; on the dual it is $-{}^tN$. The following
three conditions are equivalent.
\begin{enumerate}
\item On $\operatorname{gr}_j^M V$, every eigenvalue of $F$
has absolute value $q^{(w+j)/2}$.
\item For each $1\leq m\leq d$, all eigenvalues on
$\ker(N:\bigwedge^m V\to\bigwedge^m V)$ have absolute value
at most $q^{mw/2}$, and all eigenvalues on
$\ker(N:\bigwedge^m V^\vee\to\bigwedge^m V^\vee)$ have absolute
value at most $q^{-mw/2}$.
\item For every $1\leq m\leq d$ there are constants
$C_m,C'_m>0$ and $A_m,A'_m\geq0$ such that, for every
positive integer $r$,
\[
 \left|\operatorname{Tr}\left(F^r\mid
             \ker N_{\wedge^m V}\right)\right|
       \leq C_m r^{A_m}q^{rmw/2},
 \quad
 \left|\operatorname{Tr}\left(F^r\mid
             \ker N_{\wedge^m V^\vee}\right)\right|
       \leq C'_m r^{A'_m}q^{-rmw/2}.
\]
\end{enumerate}
When these conditions hold, condition (3) holds with $A_m=A'_m=0$
and the respective dimensions in place of $C_m,C'_m$ (or any
positive constants when a space is zero).

\begin{proof}
We first justify replacing $F$ by its semisimple part $F_s$.
If $V_\lambda$ is the generalized $\lambda$-eigenspace of $F$,
the covariance identity implies
$N(V_\lambda)\subset V_{\lambda/q}$.
Since $F_s$ is scalar $\lambda$ on $V_\lambda$, it follows
that $F_sNF_s^{-1}=q^{-1}N$. The unipotent part
$F_u=F_s^{-1}F$ commutes with $N$ and $F_s$.
The same statements hold on exterior powers and duals.
Their kernels of $N$ and all monodromy-graded spaces are
preserved by both parts. On any such subquotient, multiplying
the semisimple operator by its commuting unipotent part
does not change its eigenvalues or the traces of its powers.
Thus all three conditions are unchanged by this replacement.

There is a decomposition into $F_s,N$-stable Jordan strings.
Here is a direct proof, including the invariant complement.
Let $l$ be the largest Jordan length of $N$ on a nonzero
remaining summand. Choose an $F_s$-eigenvector $v$ of
eigenvalue $\mu$ with $N^{l-1}v\ne0$, and put
$e_j=N^jv$ for $0\leq j<l$.
Then $F_se_j=\mu q^{-j}e_j$. Choose an eigenfunctional
$a$ of $F_s$ with
$a(F_sx)=\mu q^{-(l-1)}a(x)$ and $a(e_{l-1})=1$.
Such a functional exists by semisimplicity. It vanishes
on $e_j$ for $j<l-1$, because those eigenvalues are distinct.
The map
\[
 P(x)=\sum_{j=0}^{l-1}a(N^{l-1-j}x)e_j
\]
is the identity on the span $W$ of the $e_j$.
The equality $N^l=0$ shows $PN=NP$, including the endpoint
terms. Also $N^rF_s=q^rF_sN^r$, so the $j$th coefficient
of $P(F_sx)$ is $\mu q^{-j}$ times that of $P(x)$.
Therefore $PF_s=F_sP$. Its kernel is an invariant
complement. Induction on dimension gives the asserted
decomposition.

For a string of length $l$, assign
$h(e_j)=l-1-2j$. This is its monodromy grading:
$N$ lowers it by two, and $N^a$ is an isomorphism from
grade $a$ to grade $-a$ whenever those grades occur.
Thus condition (1) on this string says
\[
       |\mu|=q^{(w+l-1)/2}.                         \tag{WM.1}
\]
The monodromy filtration of the direct sum is the sum
of these filtrations, by its defining properties.

Assume condition (2). The determinant line
$\bigwedge^lW$ is contained in $\ker N$ in
$\bigwedge^lV$, and its eigenvalue is
$\delta=\mu^l q^{-l(l-1)/2}$.
The dual determinant line is an invariant line in
$\bigwedge^lV^\vee$, because the direct-sum decomposition
also splits the dual representation. The two bounds with
$m=l\leq d$ give
$|\delta|\leq q^{lw/2}$ and
$|\delta|^{-1}\leq q^{-lw/2}$, hence equality.
Solving it gives (WM.1) for every string, proving (1).

Assume condition (1). On each string define
\[
 H e_j=(l-1-2j)e_j,\qquad
 E e_j=j(l-j)e_{j-1},\quad E e_0=0.
\]
Direct multiplication gives $[E,N]=H$.
Make the strings orthogonal, and choose their basis norms
recursively by
\[
 \|e_{j+1}\|^2=(j+1)(l-j-1)\|e_j\|^2
       \quad(0\leq j<l-1).
\]
All these factors are positive, and these norms give
$N^*=E$. On $\bigwedge^mV$ with its induced positive
Hermitian form, the induced derivations again satisfy
$N^*=E$ and $[E,N]=H$.
If $x$ is in $\ker N$ and has $H$-weight $h$, then
\[
 h\|x\|^2=\langle [E,N]x,x\rangle=-\|Ex\|^2\leq0.
\]
Thus $h\leq0$ for every nonzero such vector.
The operators $F_s$ and $H$ commute and are semisimple.
By (WM.1), on $\bigwedge^mV$ each simultaneous
$F_s,H$-eigenvector of $H$-weight $h$ has Frobenius
absolute value $q^{(mw+h)/2}$: multiply the eigenvalue
absolute values of the chosen wedge basis.
The kernel of $N$ is stable under $F_s$ and $H$, so it
has a simultaneous eigenbasis. Its eigenvalues therefore
obey the first bound of (2). The dual strings have the
same lengths and central weight $-w$. Applying the
same argument to them gives the second bound.

Condition (2) implies (3), with no power of $r$,
by summing the eigenvalues of $F^r$ with their
algebraic multiplicities. For the converse we give
the required elementary power-sum argument. For a
finite-dimensional invertible operator with distinct
eigenvalues $\alpha$ and multiplicities $b_\alpha>0$,
\[
 \sum_{r\geq1}\operatorname{Tr}(F^r)z^r
       =\sum_\alpha\frac{b_\alpha\alpha z}{1-\alpha z}.
\]
A bound $|\operatorname{Tr}(F^r)|\leq Cr^AR^r$
makes the left side holomorphic for $|z|<R^{-1}$.
If $|\alpha|>R$, the right side has a pole at
$z=\alpha^{-1}$ inside that disc. The other terms
are holomorphic there, and $b_\alpha\ne0$, so the
pole cannot cancel. The identity initially holds
in a smaller disc and hence contradicts analytic
continuation. Therefore $|\alpha|\leq R$.
Apply this to the two invariant spaces in (3),
with $R=q^{mw/2}$ and $R=q^{-mw/2}$ respectively.
Zero-dimensional invariant spaces require no argument.
This proves (2) and completes all implications.
\end{proof}

\paragraph{Geometric use and the remaining gate.}
For an $\ell$-adic representation with unipotent inertia,
inertia invariants equal $\ker N$. Apply BR-WM-02
after every complex embedding of the algebraic Frobenius
eigenvalues. For $V=H^i(X_{\overline K},\mathbb Q_\ell)$,
take $w=i$ and $d=\dim V$. Weight--monodromy for this
specified $V$ is therefore equivalent to the two
families of trace bounds in (3), for only
$1\leq m\leq d$. A finite extension of $K$ first
makes the inertia unipotent; the weight assertion
is invariant under that extension.

This is a strictly finite list of representations,
although the trace bounds still quantify over all
positive Frobenius powers $r$. It uses actual
inertia invariants and both dual directions. Bounds
only for alternating total traces, only for
$\det(1-TF\mid V)$, or only in one dual direction
are not substituted for these hypotheses.
In a geometric realization, exterior projectors
place these representations inside tensor powers
of the cohomology, so a prospective positivity
argument has explicit targets. The adopted
Dirichlet zero-location theorem alone does not
establish these local trace bounds.
